\section{Overview}\label{sec:overview}

\autoref{sec:intro} motivated over- versus underapproximate refinements, their
interaction with compositional reachability reasoning, and the extra challenges
raised by alternating modalities along curried arrows.
Here we specialize to the standard consequence readings for primitive types and,
for arrows, summarize representative modality combinations used throughout the
paper (\autoref{tab:four-modal-combinations}).
We then revisit classic refinement and coverage instantiations briefly before
examining naive mixtures and motivating strategic contexts (\emph{disjointness}).

Writing $\mstep{e}{v}$ when a term $e$ (non-deterministically) terminates at value $v$
in a call-by-value small-step semantics, over- and underapproximate base (primitive) types are often explained as follows:
\begin{align*}
    \vdash e : \ort{b}{\phi} &\iff \forall v.\; \mstep{e}{v} \implies \phi[\vnu \mapsto v]
    \\\vdash e : \urt{b}{\phi} &\iff \forall v.\; \phi[\vnu \mapsto v] \implies \mstep{e}{v}
\end{align*}
The reversal of the implication is the crux of the distinction.
If we write the behaviors of $e$ as $\{v \mid \mstep{e}{v}\}$, then in
$\ort{b}{\phi}$ the logical qualifier $\phi$ describes a superset (upper bound) of
that set, whereas in $\urt{b}{\phi}$ it describes a subset (lower bound).

\paragraph{Existing functional instances.}
With these over- and under-approximation modalities, existing work reasons about programs in functional languages.
The traditional refinement types give an overapproximate interpretation to both parameter and return positions. For example, the type $\ort{\Int}{\vnu > 0} \sarr \ort{\Int}{\vnu > 0}$ indicates a function may only return a positive integer (or diverge) when the input argument is also positive. This is analogous to Hoare triples $\{P\}\;e \;\{Q\}$ whose pre- and postconditions are both overapproximate and thereby establish safety:
``on every possible execution, the program does not go wrong.''
Recent \emph{coverage types}~\cite{CoverageType} instead interpret the
\emph{return} type of a function underapproximately, which supports reasoning about
\emph{coverage completeness} of input generators. For example, the type $x{:}\ort{\Int}{\vnu \geq 0} \sarr \urt{\List[\Int]}{\I{len}(\vnu) = x}$ indicates a sized list generator that must produce all lists of length $x$. 
The combination of overapproximate parameter and underapproximate return provides a way to do must-style reasoning and yields a lower-bound guarantee on the outputs of a nondeterministic program.

\subsection{Under- and Over-approximation Typing for Functional Programs}

A natural next step is to explore the remaining combinations of modalities
in function types.
Coverage types forbid an underapproximate interpretation of the parameter
position.
Yet one can imagine pairing an underapproximate parameter with an
underapproximate return, much like an incorrectness triple $[P]\;e\;[Q]$, to certify
\emph{incorrectness}-style properties of functions.
For example, a function $f$ of type
$\urt{\Int}{\top} \sarr \urt{\Int}{\vnu = 0}$ would mean that there exists
some integer input on which $f$ \emph{must} (terminatingly) return $0$,
witnessing a potential division-by-zero exception such as $3 / (f\;x)$ (if we assume $x$ can cover all integers).
Dually, under parameter plus over return is weaker. For example, the type
$\urt{\Int}{\top} \sarr \ort{\Int}{\vnu = 0}$
does not guarantee termination, so one cannot in general prove that a potential error \emph{must} occur.
However, it is close to bug detection in program analysis, where false positives are traded for breadth of alarms. Many classical static analyses exploit this under parameter plus over return combination effectively.

\begin{table}[ht!]
  \centering
  \caption{Representative modality combinations at function types. Meanings correspond to consequence readings spelled out later in~\autoref{sec:formal}; strategy arrows \(\ctsarr\) make argument dependence syntactic; sums \(\tau_1 \ctoplus \tau_2\) keep branching visible at types.}
  \label{tab:four-modal-combinations}
  \scriptsize
  \setlength{\tabcolsep}{5pt}
  \begin{tabular}{@{}>{\raggedright\arraybackslash}p{0.29\linewidth}>{\raggedright\arraybackslash}p{0.36\linewidth}>{\raggedright\arraybackslash}p{0.31\linewidth}@{}}
    \toprule
    \textbf{Typing shape for \(f\)} & \textbf{Meaning (consequence form)} & \textbf{Typical use} \\
    \midrule
    $\emptyset \ctvdash f : x{:}\ort{b_x}{\phi_x} \ctsarr \ort{b}{\phi}$ &
    If $x$ satisfies $\phi_x$ and $f\;x$ terminates at $v$, then $v$ satisfies $\phi$. &
    Refinement-type safety (over--over), Hoare-style obligations. \\
    $\emptyset \ctvdash f : x{:}\ort{b_x}{\phi_x} \ctsarr \urt{b}{\phi}$ &
    If $x$ satisfies $\phi_x$ and output $y$ satisfies $\phi$, then $f\;x$ can reach \(y\) (coverage on returns). &
    Coverage / generator completeness~\cite{CoverageType}. \\
    $\emptyset \ctvdash f : x{:}\urt{b_x}{\phi_x} \ctsarr \urt{b}{\phi}$ &
    For each output \(y\) with \(\phi[\vnu\mapsto y]\), there is some input \(x\) in \(\phi_x\) reaching \(y\) (incorrectness \(\forall_{\mathrm{post}}\exists_{\mathrm{pre}}\)). &
    Incorrectness triples~\cite{OHP19}. \\
    $\emptyset \ctvdash f : x{:}\urt{b_x}{\phi_x} \ctsarr \ort{b}{\phi}$ &
    There exist \(x\) in \(\phi_x\) and output \(y\) with \(\phi[\vnu\mapsto y]\) reached by \(f\,x\) (existential linkage).&
    Existential summaries; static bug reports that pick one witness. \\
    $\emptyset \ctvdash f : x{:}\ort{b_x}{\phi_x} \ctsarr \ort{b}{\phi} \ctoplus b$ &
    For every \(x\) in \(\phi_x\), some branch of \(f\) yields an output satisfying \(\phi\) (sufficient incorrectness with control-flow split). &
    Must-reach arguments under over-bounded parameters. \\
    $\emptyset \ctvdash f : x{:}\urt{b_x}{\phi_x} \ctsarr \ort{b}{\phi} \ctoplus b$ &
    Existential witness where some \(\ctoplus\)-summand reaches \(\phi\) (combined with erased base \(b\) on sibling branches). &
    Branch-local witnesses under under-bounded inputs. \\
    \bottomrule
  \end{tabular}
\end{table}

\autoref{tab:four-modal-combinations} organizes the modality patterns we use in later rules; branching on return types~\(\ctoplus\) matches the \(\ctoplus\)-lifting in~\autoref{sec:formal} for ``OR'' summaries without collapsing independent obligations.
Allowing each position to swing between modalities is tempting but runs into theoretical difficulties.
Simply dropping coverage types' restriction that parameters cannot be
underapproximate leads to multiple problems.
Some terms behave almost like \emph{substructurally constrained} resources.
Consider an identity function $\Code{f_{id}} \doteq \lambda x.\,x$.
It admits an incorrectness-flavoured specification
$\urt{\Int}{\top} \sarr \urt{\Int}{\top}$: for every possible result $y$ of
an application, there is an input $x$ with $y = f\;x$ (take $x = y$).
Yet consider the program
\begin{minted}[xleftmargin=5pt, numbersep=4pt, linenos = true, fontsize = \small, escapeinside=??]{OCaml}
    let y = ?$\Code{f_{id}}$? (int_gen ()) in
    let z = ?$\Code{f_{id}}$? (int_gen ()) in
    y - z
\end{minted}
Individually, $y$ and $z$ might have type
$\urt{\Int}{\top}$ which means that they can cover all integers. As a consequence, $y - z$ should also cover all integers (e.g., let $z = 0$ and impose no constraint on $y$). However,
$y - z$ can only be $0$.
Although $y$ and $z$ separately have coverage guarantees similar to a random generator like \Code{int\_gen}, they are always equal.
The two outcomes are \emph{correlated} rather than independently drawable under witnesses (we borrow probabilistic language, but there is no probability here); our calculus records that correlation in the bunching of the context rather than by ``spending'' a witness on successful termination.

The issues of correlation among return values are not the whole story.
A call can also change which refinements may soundly appear \emph{together} on variables that remain in scope---not because our core language crashes and thereby deletes executions, but because the continuation must respect \emph{dependencies} introduced by the call.
Consider the following \emph{total} example (our core calculus omits failure primitives such as \Code{err}):
\begin{minted}[xleftmargin=5pt, numbersep=4pt, linenos = true, fontsize = \small, escapeinside=??]{OCaml}
  let f_abs x = if x >= 0 then x else -x in
  let x = int_gen () in
  let y = ?$\Code{f_{abs}}$? x in
  x
\end{minted}
The function $\Code{f_{abs}}$ has type
$\urt{\Int}{\top} \sarr \urt{\Int}{\top}$: for every nonnegative result $y$ there is some input $x$ with $|x| = y$.
After the call, $x$ is still the same integer drawn from \Code{int\_gen}, but any reasoning that treats $x$ and $y$ as two $\sep$-independent under-complete draws would overstate what is known, because the total definition forces $y = |x|$.
This is a phenomenon of \emph{additive} bunching and dependency, not of ruling out classes of executions through unmodeled failure (the core language has no crash primitive).

A third issue is higher-order behavior:
\begin{minted}[xleftmargin=5pt, numbersep=4pt, linenos = true, fontsize = \small, escapeinside=??]{OCaml}
  let f (g: unit -> bool) : bool =
    let y = g () in
    let z = g () in
    y == z in
  let g' = (fun () -> true) ?$\oplus$? (fun () -> false) in
  f g'
\end{minted}
The function parameter $g$ has type $\Unit \sarr \urt{\Bool}{\top}$. Note that the parameter type happens to be $\Unit$, which is the same under both over- and under-approximation, so ``by rights'' $y$ and $z$ ought to be independent draws.
One might expect $y$ and
$z$ each to enjoy $\urt{\Bool}{\top}$, where both $\true$ and $\false$ are reachable.
Yet the actual argument $g'$ is implemented as a nondeterministic choice ($\oplus$) between
two functions that return only \texttt{true} or only \texttt{false}.
This term $g'$ can still be given type $\Unit \sarr \urt{\Bool}{\top}$; both $\true$ and $\false$ are reachable.
However, $y$ and $z$ in the higher-order function application can only be both \texttt{true} or both \texttt{false} in a call-by-value semantics.
On the other hand, if $g$ were used only once inside $f$, there would be no such issue. For example, if $f$
returned just $y$, one could still ascribe $\urt{\Bool}{\top}$ to
$f\;g'$.

In short, underapproximate parameters capture an attractive \emph{external user-specified verification goal}, i.e., ``there exists an input''.
What the bare modalities omit is how a call reshapes whatever else remains live in scope: correlations between successive results, weakening of leftover argument guarantees, or higher-order reuse that forces shared outcomes.
Those interaction questions are all about \emph{disjointness} between strategy resources, not about the arrow $\sarr$ alone.

\subsection{Logic for Over- and Under-approximation Function Typing}

These examples show that freely mixing under- and overapproximation in
function types is problematic and can yield unsound reasoning.
One might patch the type system with linear constraints or substructural typing rules, but
that may lose expressiveness on reachability reasoning.
What we really need is a more accurate \emph{logical} interpretation of function
application than off-the-shelf correctness or incorrectness logics provide when their native unit is a global store.
(A brief comparison with outcome logic's branching idiom is deferred to~\autoref{sec:related}.)

Consider a functional program
$\mathtt{let}\; y = f\;x\;\mathtt{in}\; \cdots$ which starts with a function application.
In a state-based logic, both $x$ and $y$ are variables in the global store, thus it is natural to require both pre- and postconditions to constrain both $x$ and $y$.
It becomes less natural when the surface language is functional and the
abstraction we actually manipulate is a \emph{function type} $\tau_x
\sarr \tau_y$: programmers think they are specifying ``what $f$
returns,'' not ``how the entire context changes after
application.''
The function in a (pure) functional language is treated as a simple abstraction that is independent from the context. Recording the context change after function application conflicts with the functional programming paradigm.

\paragraph{Concrete illustrations.}
Take $\Code{f_{id}} \doteq \lambda x.\,x$ as a running example.
The most informative incorrectness triple for a single call is
\[
  [x > 0]\quad y \mathrel{:=} \Code{f_{id}}\;x \quad [y = x \land x > 0],
\]
where the postcondition $y = x \land x > 0$ explicitly records the input-output relationship.
Writing a looser postcondition, for example $[y > 0]$, keeps the same lower bound guarantee on $y$ but loses the relationship between $x$ and $y$, which makes the triple unsound. For example, the state $[x\mapsto 1, y\mapsto 2]$ is allowed by the simplified postcondition but is not reachable by $\Code{f_{id}}$. This issue is already discussed in \cite{OHP19}: every underapproximate postcondition needs to be precise so that it allows only reachable post-states. This requirement is too strong for functions in our setting.

% The argument-side is symmetric.
% $\Code{f_{pos}} \doteq \lambda x.\,
% \zif{x > 0}{\Code{int\_gen()}}{\Code{err}}$
% admits an informative IL triple
% \[
%   [x > 0]\quad y \mathrel{:=} \Code{f_{pos}}\;x
%   \quad [y \in \mathbb{Z} \wedge x > 0],
% \]
% where $x > 0$ must be \emph{re-stated} in $Q$ lest the constraint on the
% argument be lost for any continuation that reads $x$.
% In a type-based account, one would instead expect the function arrow
% itself, something like $\urt{\Int}{x > 0} \sarr \urt{\Int}{\top}$, to
% carry this information structurally, without forcing the caller to repeat it
% in every surrounding assertion.

The examples above show that lifting imperative program logics verbatim pushes to the opposite extreme: the underapproximate return
type must be precise, which means it must be correlated with the parameter.
That erodes much of what makes refinement types pleasant.
For instance, when approximation modalities compose freely, an
overapproximate return need not mention the parameter, whereas an
underapproximate one seemingly must. This is an asymmetry.
Moreover, not all parameters need to be correlated with the return. Coverage types already show that an overapproximate parameter can be independent from the return.
Worse, forcing everything through a monolithic postcondition breaks
\emph{currying}: the final return type is special, so
$x{:}\tau_x \sarr y{:}\tau_y \sarr \tau$ cannot be rearranged into
$x{:}\tau_x \sarr (y{:}\tau_y \sarr \tau)$ without changing the meaning of
the return position.

We therefore aim for a new logic that does \emph{not} require fully
explicit expressions of the precise relation between function parameter and return, but still distinguishes the cases that
matter: whether a function may be invoked more than once, how parameters evolve under calls, and whether multiple outputs are
independent.
These are the questions that resource- or permission-based logics (e.g.\
separation logic) were built to address: output independence as a form of
disjointness, parameter evolution as permission separation (read versus
write, metaphorically), and higher-order reuse as a question of how strategic information is shared between invocations.

Importantly, we are \emph{not} importing the usual operational reading of
``resources'' as memory or explicit effects, nor do we claim a probabilistic
notion of correlation.
Our resources arise from the \emph{proof-relevant} distinction between
angelic and demonic choice in program reasoning. That structure is imposed by the
logic, not by the underlying operational semantics.
Even for a deterministic, heap-free program such as
$\Code{f_{id}}$, this form of resource reasoning is both available
and necessary.

\subsection{Strategy Logic and Strategy Types}

We develop \emph{strategy logic}---a resource-based modal logic---and matching \emph{strategy types}, in which over- and underapproximation are explicit modalities. As motivated above, the framework need not force fully explicit descriptions of the exact dependence of returns on parameters, yet it still distinguishes the cases that matter. In the same spirit as Hoare logic, which reasons about a program under multiple possible execution contexts (for instance, requiring that every state satisfying the precondition reach a state satisfying the postcondition), judgments are relativized to proof-relevant \emph{strategic contexts}: each hypothesis carries strategy resources that record which assignments remain simultaneously viable. That packaging lets us express disjointness of strategic fragments together with read versus write permission---and finer structure at function types (sums \(\ctoplus\), \(\ctsarr\) versus \(\wand\))---without insisting on imperative-style postconditions over a single global store.
As in~\autoref{sec:formal}, strategy-type hypotheses are \emph{not} consumed by successful termination; typing tracks \emph{who may depend on whom} (and when an ambient coercion such as \(\Code{ToOver}\) is needed), rather than decrementing a witness budget on crashes.

\paragraph{Strategy types as resources} The core idea is to read an underapproximate strategy type as a resource that is not duplicable. For example,
\begin{align*}
  &x{:}\urt{\Int}{\vnu > 0} \vdash x : \urt{\Int}{\vnu = 1} \quad \judgok
  \\&x{:}\urt{\Int}{\vnu > 0} \vdash x : \urt{\Int}{\vnu = 2} \quad \judgok
  \\&x{:}\urt{\Int}{\vnu > 0} \vdash (x, x) : \urt{\Int \times \Int}{\vnu.1 = 1 \land \vnu.2 = 2} \quad \judgfail
\end{align*}
Here $x$ is annotated in the typing context with an underapproximation type $\urt{\Int}{\vnu > 0}$, indicating writable strategy information for $x$: refinements may still ``write'' or steer witnesses within the positive integers as needed. The first two lines therefore show that $x$ can be proved equal to $1$ and to $2$, respectively. The last line shows that the pair $(x, x)$ cannot be proved equal to $(1, 2)$: one cannot treat the two components as \emph{independent} witnesses for separate positive integers, because they are the same $x$ (contraction is disallowed in~\autoref{sec:formal} for exactly this reason). Overapproximation, by contrast, behaves like a resource with ``read'' permission:
\begin{align*}
  &x{:}\ort{\Int}{\vnu > 0} \vdash x : \ort{\Int}{\vnu > -1} \quad \judgok
  \\&x{:}\ort{\Int}{\vnu > 0} \vdash x : \ort{\Int}{\vnu > -2} \quad \judgok
  \\&x{:}\ort{\Int}{\vnu > 0} \vdash (x, x) : \ort{\Int \times \Int}{\vnu.1 > -1 \land \vnu.2 > -2} \quad \judgok
\end{align*}
Again $x$ carries an overapproximation type $\ort{\Int}{\vnu > 0}$: the value of $x$ is \emph{decided by others} and treated as ``read-only''. If $x$ is read as some positive integer, it may be shown strictly greater than $-1$ and than $-2$, respectively. Read permission may be duplicated, so both components of $(x, x)$ may simultaneously carry these constraints, unlike under the underapproximation reading.

\paragraph{Disjointness} Separation logic uses $\sep$ to express disjointness of heap resources; analogously, we express disjointness of strategic contexts (equivalently: independence of the strategy resources backing them).
Contexts use three related connectives whose differences matter as soon as more than one under-refinement stays live:
\begin{itemize}
\item \(\Gamma_1 \sep \Gamma_2\) asserts \emph{full disjointness}: the two halves are combined multiplicatively, so neither binds variables carrying strategy fragments that overlap the other's unless those fragments genuinely split (the usual separation-logic intuition).
\item \(\Gamma_1 , \Gamma_2\) plus an \emph{explicit} dependency arises when hypotheses on the right mention variables from the left: the second refinement may refer directly to witnesses bound earlier.
\item \(\Gamma_1 , \Gamma_2\) with no cross-variable occurrences can still impose \emph{implicit} sequential dependence---the additive bunch is checked after the ambient context absorbs the left block, recording possible correlation even without surface-level name sharing.
\end{itemize}
\noindent
For illustration, juxtaposition without separation rejects independent pairings whenever \(y\) could have been instantiated from \(x\) (e.g.\ \(\zlet{y}{x}{\ldots}\)), whereas separating restores genuine independence when two draws are disjoint:
\begin{align*}
  &x{:}\urt{\Int}{\vnu > 0}, y{:}\urt{\Int}{\vnu > 0} \vdash (x, y) : \urt{\Int \times \Int}{\vnu.1 = 1 \land \vnu.2 = 2} \quad \judgfail
  \\&x{:}\urt{\Int}{\vnu > 0} * y{:}\urt{\Int}{\vnu > 0} \vdash (x, y) : \urt{\Int \times \Int}{\vnu.1 = 1 \land \vnu.2 = 2} \quad \judgok
\end{align*}

\paragraph{Functional parameters versus outer context.}
\emph{(Example forthcoming.)}
This pattern asks whether introducing (or refining) an underapproximate formal parameter clashes with disjointness assumptions on the enclosing context---for instance whether the caller's strategic context may still be split orthogonally once the callee's witness is unpacked.

\paragraph{Call results versus outer context.}
\emph{(Example forthcoming.)}
Symmetrically, this pattern asks whether a freshly bound return behaves as an \(\sep\)-independent observer or must stay correlated with lingering assumptions after \(f\,x\)---the situation that distinguishes ``two draws'' from ``same coin twice.''

\paragraph{Function application} Reading over- and underapproximation as read and write permission on strategy resources, and with disjointness in hand, we can disentangle the situations that arise from function application without precisely tracking the functional relation between argument and result.
Consider, for instance,
\begin{align*}
  &x{:}\urt{\Int}{\top} \vdash \zlet{y}{f\;x}{e} : \tau
\end{align*}
and focus on the \emph{context} in which the continuation $e$ is typed after the call. For simplicity, suppose $f\;x$ can cover all integers; the following three cases must then be separated:
\begin{enumerate}
\item The continuation may keep $x$ \emph{over-bounded} (read-only) while still learning an under-complete return $y$---as when \textsc{T-LetD} applies an ambient $\Code{ToOver}$ coercion in~\autoref{sec:formal}, rather than because execution failed and deleted a witness.
\item $x$ and $y$ both stay under-complete, but the bunch governing them is \emph{correlated} (additive). For example, let $f$ be $\Code{f_{id}}$ or the total $\Code{f_{abs}}$ from above.
\item $x$ and $y$ are \emph{independent} under-complete draws, separated by $\sep$. For example, let $f$ be $\lambda x.\Code{int\_gen()}$, which returns a nondeterministically chosen integer unrelated to the witness for~$x$.
\end{enumerate}

\paragraph{Separating \texttt{let} and \textsc{T-LetD}.}
Rule \textsc{T-LetD} in~\autoref{sec:formal} types a \texttt{let}-binder whose continuation is checked under a separating extension $\Gamma \ctstar y{:}\tau$, while the right-hand side is judged in a related ambient~$\Gamma'$ with $\Code{ToOver}(\Gamma,\Gamma')$.
Concretely, after an under-complete $x$ is introduced, a second draw may be required $\sep$-independent of the witness for~$x$:
\begin{minted}[xleftmargin=5pt, numbersep=4pt, linenos = false, fontsize = \small, escapeinside=??]{OCaml}
  let x = int_gen () in
  let y = int_gen () in
  (* continuation uses both x and y *)
\end{minted}
\noindent
For the inner \texttt{let}, \Code{int\_gen()} is typed in the coerced context~$\Gamma'$ so that outer variables such as $x$ are read over-bounded for that premise alone; the continuation is then checked with $y$ $\sep$-appended to the outer bunch, yielding the schematic shape of case~1 above ($x$ treated read-only alongside an under-complete~$y$).
This is the same $\Code{ToOver}$ discipline highlighted in the caption of Figure~\ref{fig:auxiliary-rules}; it records \emph{dependency control} for separating binds, not consumption of coverage by failure.

These regimes correspond to the following typing contexts:
\begin{align*}
  &x{:}\ort{\Int}{\top}, y{:}\urt{\Int}{\top} \vdash e : \tau \quad \text{(case 1)}
  \\&x{:}\urt{\Int}{\top}, y{:}\urt{\Int}{\top} \vdash e : \tau \quad \text{(case 2)}
  \\&x{:}\urt{\Int}{\top} * y{:}\urt{\Int}{\top} \vdash e : \tau \quad \text{(case 3)}
\end{align*}
Function types for $f$ that induce these three contexts can be expressed using the ordinary arrow $\sarr$ and the magic wand $\wand$; the next section develops the details.

\paragraph{Higher-order functions and Persistence}
Different from the first order case, functional language that supports higher-order functions, where the parameter of a function can also be another function. This extension introduce more chelleges into our modal type system. Consider the following type context:
\begin{align*}
  &x{:}\ort{\Int}{\top} \vdash (f\;x) + (g\;x) : \tau  \quad \judgok \quad &\text{(when $f$ and $g$ has type $\ort{\Int}{\top} \sarr \urt{\Int}{\vnu = 1}$)}
  \\&x{:}\urt{\Int}{\top} \vdash (f\;x) + (g\;x) : \tau \quad \judgfail \quad &\text{(when $f$ and $g$ has type $\urt{\Int}{\top} \sarr \urt{\Int}{\vnu = 1}$)}
\end{align*}
Unlike overapproximation, the variable $x$ with underapproximation type cannot be used for multiple times in the second line, since the $f$ and $g$ has hidden and potentially different assumptions about their parameters. In a first order setting, the overapproximation and underapproximation is sufficient to distinguish weather a variable can be used for multiple times.
However, when we have function parameters, the thing goes more complicated since a function type is allowed to mixed these two modalities in our setting. Moreover, even for the same function types, the case that allowed to be used for multiple times or not are still coexisting. Consider the following example:
\begin{align*}
  &x{:}\urt{\Bool}{\top} \vdash f_1 \doteq \lambda ().(\true \oplus \false) : \Unit\sarr\urt{\Bool}{\top} \quad \judgok
  \\&x{:}\urt{\Bool}{\top} \vdash f_2 \doteq \lambda ().\true \oplus \lambda ().\false : \Unit\sarr\urt{\Bool}{\top} \quad \judgok 
  \\&x{:}\urt{\Bool}{\top} \vdash f_3 \doteq \lambda ().x : \Unit\sarr\urt{\Bool}{\top} \quad \judgok
\end{align*}\noindent
The three functions $f_1$, $f_2$, and $f_3$ above both has type $\Unit\sarr\urt{\Bool}{\top}$, which means that they must non-deterministically returns all boolean values. The difference is that $f_1$ encapsulates the non-determinism into its function body, $f_2$ exposes the non-determinism to top level, where $f_3$ is depends on non-determinitsm from context (i.e., $\zlet{x}{\Code{int\_gen()}}{f_3}$). These three functions looks similar, while only $f_1$ can be used for multiple times. Consider the following use case 
\begin{align*}
  &x{:}\urt{\Bool}{\top} \vdash g \doteq \lambda f.(f\;() == f\;()) : f{:}(\Unit\sarr\urt{\Bool}{\top})\sarr\urt{\Bool}{\vnu = \false} \quad \textbf{\textcolor{red}{???}}
\end{align*}
The higher-order function $g$ takes a function and applies it for twices, and expects the result can be different. We show the reduction results of apply $g$ for $f_1$, $f_2$, and $f_3$ respectively:
\begin{align*}
  g\;f_1 &\hookrightarrow^* \true == \true \text{  or  } \true == \false \text{  or  } \false == \true \text{  or  } \false == \false \quad \judgok
  \\&\hookrightarrow^* \true \text{  or  } \false \quad \judgfail
  \\g\;f_2 &\hookrightarrow^* \true == \true \text{  or  } \false == \false \hookrightarrow^* \true 
  \\g\;f_3 &\hookrightarrow^* x == x \hookrightarrow^* \true \quad \judgfail
\end{align*}
As shown above, only $g\;f_1$ can result to the value $\false$, and two other cases can only result to the value $\true$. The intuition here is that only the non-determinism that is encapsulated in function body are duplicatable, since it is independent from the decision of external context of operational semantics.

In order to express this decision independence, we introduce a the \emph{persistency} modality $\ctpers$ as other resource-based logics and types systems do. The key intuition is that the \emph{singleton} resource (i.e., the resource contains only one store) remains no decision to make, thus a type judgement holds under singleton resource is persistent. We would highlight that the resource we consider is both derived from the context and also the decision made by the operational semantics. For example, the context relied by $f_3$ is not singleton, and the term $f_2$ will not just reduce to a single value, thus both of them cannot have persistent type. On the other hand, the parameter of $g$ is required to be persistent. In general, the overapproximation type in the type context can be treated as persistent, while typically the underapproximation type is not persistent, and function types should be explictly marked as persistent or not.
In order to introduce a persistent type in our type system (i.e., $\Gamma \ctvdash e : \tau$), we require the type context $\Gamma$ only consist with persistent types (to rule out the case like $f_3$) and $e$ itself should be single valued (to rule out the case like $f_2$).
For example, the following type judgement holds:
\begin{align*}
  &x{:}\ort{\Int}{\vnu > 0} \vdash \lambda ().\zif{x > 0}{(\true \oplus \false)}{\true} : \ctpers (\Unit\sarr\urt{\Bool}{\top}) \quad \judgok
\end{align*}\noindent
On the other hand, the parameter type of $g$ should also marked as persistent, since input function parameter used twice in $g$.

% Under-complete refinements behave like \emph{non-duplicable} witnesses in the sense of~\autoref{sec:formal} (no naive contraction), whereas overapproximate refinements behave read-only---and function arrows themselves behave like persistent specifications in ordinary refinement systems.
% Higher-order nondeterminism nevertheless forces a sharper distinction between \emph{where} branching lives.
% Consider
% \begin{align*}
%   &g_1 \doteq (\lambda ().\true) \oplus (\lambda ().\false)
%   \\&g_2 \doteq \lambda ().(\true \oplus \false).
% \end{align*}
% After one application either term can expose both Boolean witnesses, but repeating the call aligns with~\autoref{sec:formal}'s discussion: collapsing the choice \emph{between} thunk bodies into a lone return refinement erases branching that must stay observable for sound reuse.
% Informally \(g_2\) admits the customary functional reading \(\Unit \sarr \urt{\Bool}{\top}\) whose internal \(\oplus\) does not correlate repeated answers the way \(g_1\) correlates \(\true\cdot\true\) with \(\false\cdot\false\).
% For \(g_1\) the type discipline instead keeps \(\oplus\) at \emph{function} granularity, exposing a \(\ctoplus\) of two unary arrows (\(\Unit\sarr\cdot\,\ctoplus\,\Unit\sarr\cdot\) as in~\autoref{sec:formal}), which recovers precisely the observable split between \(\lambda ().\mathit{b}_1\) and \(\lambda ().\mathit{b}_2\) callers must respect.

% \paragraph{Strategy logic and strategy types} ...

% \paragraph{Resources as possible executions, not states.}
% We therefore move from predicates on a global store to \emph{resources}:
% finite collections of \emph{partial substitutions}, each representing a
% consistent way the program might resolve nondeterministic or unknown
% choices.
% Combining resources uses a separating product (compatible glueing of
% partial information) \emph{and} a strategy-branching sum (branching), so branching,
% independent correlation, and framing need not be encoded as mutations
% of an implicit heap.
% The carrier is purpose-built for ``which execution possibilities are
% we keeping track of?'' rather than ``what is the unique concrete
% state?'' Hence the semantics does not presuppose a stateful language.

% \paragraph{Strategy logic and strategy types.}
% On this algebra we define \emph{strategy logic}, a modal bunched logic
% with approximation operators that play the roles of upper and lower
% bounds with respect to choice.
% It provides the missing logical distinctions (products,
% wands, sums, and modalities) that let us express the three application
% regimes compositionally.
% \emph{Strategy types} are the layer above: refinement constructors export
% the modalities to data, while function constructors $\sarr$ and $\wand$
% reify exactly the parameter/return dependency patterns listed above,
% all interpreted by uniform clauses into strategy logic.

% \paragraph{Roadmap.}
% \autoref{sec:formal} presents substitutions, the strategy algebra, pure
% propositions embedded in it, basic and approximate strategy logic, and
% the syntax and denotation of modal refinement types.
% Syntactic typing rules are deferred; the denotational definition fixes
% the specification they should implement.
